\documentclass[11pt]{article}
\usepackage{amsmath,amssymb,amsthm}
\newtheorem{theorem}{Theorem}
\title{What a retained formation proof needs to know about later constraints}
\begin{document}
\maketitle

A local controller can reuse a proof after a change without retaining the
whole constraint graph. It must retain the bounds between locations where
later requirements can attach. Checking only the individual locations, or
checking changed requirements separately, need not preserve admission.

\paragraph{Lift and query family.}
A finite feasible base \(G\) constrains real variables by inequalities
\(x_v-x_u\le c_{uv}\). For the single-file model these are differences
of positions; for a simple temporal network they are differences of
times. For absolute choices include a reference index in the interface
with prescribed value zero. Let \(S\) be the
declared set of indices, the only endpoints of replaceable requirements.
Future changes can add any finite batch of difference inequalities
between members of \(S\). A replaceable requirement must be kept outside
the frozen base so that removing it does not leave a stale implication.
The consumer asks for exact feasibility of the whole batch, and, if
needed, an extension of a feasible assignment on \(S\). It does not ask
for physical observations, unknown route changes in \(G\), or a policy
under uncertain observations.

\begin{theorem}
For the feasibility consumer, the exact retained information is the
matrix \(d(u,v)\) of base shortest-path bounds between indices in \(S\),
with \(+\infty\) for unavailable paths. Two feasible bases have the
same matrix iff they agree on every admissible batch, even when only
single-edge tests are used to distinguish them. Every assignment on
\(S\) satisfying the retained bounds lifts to a base assignment.
\end{theorem}
\begin{proof}
Any base solution satisfies the path bounds. Conversely suppose the
declared values \(y_s\) satisfy them. Use one base feasible potential
\(q\), and choose \(M\ge\max_{s\in S}(y_s-q_s)\). Set
\[
   x_w=\min\bigl(M+q_w,\ \min_{s\in S}(y_s+d(s,w))\bigr).
\]
Ignore infinite terms, and use \(q\) alone if no values are declared.
Each finite-distance term follows the edge inequalities by the path
triangle inequality. Taking pointwise minima preserves those
inequalities; the \(q\) term covers vertices not reachable from \(S\).
At any \(s\in S\), the \(s\)-term gives \(y_s\), and no other term
can be smaller, by the retained bounds and the choice of \(M\).
This gives an extension and preserves every new constraint on \(S\).
It is an exact reduction for the declared consumer.

To distinguish unequal matrices, take entries \(d_1(u,v)<d_2(u,v)\).
Add the single reverse edge \(v\to u\) of weight \(c\).
The base is feasible, so inconsistency arises precisely if the new
edge closes a path of total weight \(d_i(u,v)+c<0\). Choose \(c\)
to make this true in the first base but not the second. This also
works if \(d_2=+\infty\), since no reverse cycle would exist there.
No exact representation of this whole query interface can collapse
the two entries.
\end{proof}

\paragraph{Size and proof work.}
With \(k\) terminals and \(q\) possible rational weights in \([1,2]\),
make a complete directed base with each of its \(k(k-1)\) off-diagonal
weights chosen independently. These are exactly the shortest-path
entries: a two-edge path costs at least two. Thus any fixed-width
summary for all admissible single-edge tests needs at least
\(\lceil k(k-1)\log_2 q\rceil\) bits in this family. The bound is
not a claim for a sparse known corridor, a fixed successor,
or a comparator given different admissible changes.

After compilation, a feasibility check reads at most \(k(k-1)\)
finite retained bounds plus the changed batch. A full interior
assignment still needs access to the base or a decoder and its
stored distances. The example compares edge counts of the receiver
check, not total proof cost, runtime, or memory including decoding.
Ordinary simple temporal network closure and incremental
propagation remain the direct comparators; this result specifies
the exact restriction and its failure mode rather than proposing
a faster general shortest-path algorithm.

\paragraph{Witness and boundary.}
An internal path implies that a formation's total span is at
most four. Two local replacements individually allow requiring
three units of separation at either gap. Installing both requires
span at least six, so the checks must share one retained outer bound.
A second case removes an old replaceable upper bound of four,
leaving base upper bound ten. Requiring span at least seven is
now feasible. A receiver that retained the old bound inside its
base still refuses; it tests the wrong law.

The exact replay compares 780 changed batches against the full graph,
lifts 371 assignments, and verifies 729 distinct three-terminal
query signatures. The chain family consumes four edges in each
compiled test while the full test grows to 501. These are defined
formation constraint systems, not independent physical missions
or full R2 acceptance.
\end{document}
